% Einheit 3 von 5 – Prozentwert berechnen (bank/prozentrechnung/e3.jsonl, zusammenbau v0.3)
%% TODO Zweigzeile Teil 1: was hier gelernt wird (Satz in Schülersprache aus dem Einheitstitel) – Einheit 3
\einheitenkopf[e3]{Einheit 3 von 5 · Prozentwert berechnen}
\zweigzeile{ab Kl. 7 · P10}
\setcounter{aufgabe}{16}

%% TODO Titel: Ich-kann-Satz fehlt in der Bank (Platzhalter: Kettenname) – Nr. 17
%% TODO Anweisung: ein Satz über der ersten Teilaufgabe fehlt in der Bank – Nr. 17
\begin{aufgabe}{Prozentwert}
%% TODO \rechenplatz{n}: Schrittzahl des Grundfalls fehlt in der Bank (nur bei fester Schreibform, 2.3 b)
\begin{teile}
\teil $50\,\%$ von $80\,€$ – wie viel?

\streifen[2]{50}{$0$}{$80\,€$}
\leerfeld[€]
\teil $6\,\%$ von $400$ kg \leerfeld[kg]
\teil $30\,\%$ von $90\,€$ \leerfeld[€]
\teil $45\,\%$ von $80$ kg – mit Dezimalzahl \leerfeld[kg]
\teil $20\,\%$ von $4{,}50\,€$ \leerfeld[€]
\teil Schuhe für $70\,€$, $20\,\%$ Rabatt – wie viel spart man? \leerfeld[€]
\teil $150\,\%$ von $40$ kg \leerfeld[kg]
\teil Ein Fernseher kostet $480\,€$. Bei Barzahlung gibt es $15\,\%$ Rabatt. Wie viel Euro spart man? Kreuze an. \hfill (P10 2021 OS)\\ \kreuz{$32\,€$}\\ \kreuz{$408\,€$}\\ \kreuz{$72\,€$}\\ \kreuz{$552\,€$}
\teil $17\,\%$ von $60\,€$ \hfill (P10 2014 OS) \\ \leerfeld[€]
\teil Bei einer Umfrage unter $800$ Haushalten hatten $35\,\%$ ein E-Bike, sechs Jahre später $52\,\%$. Um wie viele Haushalte hat die Zahl zugenommen? \hfill (P10 2019 OS)
\end{teile}
\end{aufgabe}

%% TODO Titel: Ich-kann-Satz fehlt in der Bank (Platzhalter: Kettenname) – Nr. 18
%% TODO Anweisung: ein Satz über der ersten Teilaufgabe fehlt in der Bank – Nr. 18
\begin{aufgabe}{Mehrwertsteuer in Euro}
\begin{teile}
\teil Reparatur ohne Mehrwertsteuer $300\,€$, dazu $19\,\%$ – wie viel Euro Steuer? \leerfeld[€]
\end{teile}
\end{aufgabe}

%% TODO Titel: Ich-kann-Satz fehlt in der Bank (Platzhalter: Kettenname) – Nr. 19
%% TODO Anweisung: ein Satz über der ersten Teilaufgabe fehlt in der Bank – Nr. 19
\begin{aufgabe}{Prozentwert – Fehler finden, Begründen, Darstellungswechsel, Anwendung}
\begin{teile}
\teil Jonas, $60\,\%$ von $35\,€$: \rechnung{0{,}06 \cdot 35 &= 2{,}10\,€} Fehler? Wie heißt es richtig?
\teil $1\,\%$ von $500\,€$ sind $5\,€$. Warum?
\teil Ganzes $60\,€$: $30\,\%$ im Streifen einzeichnen, Wert ablesen.

\streifen[0]{0}{$0$}{$60\,€$}
\leerfeld[€]
\teil Dieselbe Jacke: in Laden A $95\,€$ mit $20\,\%$ Rabatt, in Laden B $85\,€$ mit $10\,\%$ Rabatt. Wo ist sie billiger?
\end{teile}
\end{aufgabe}
